\section{Results}\label{sec:results}

We now apply the posteriors of Table~\ref{tab:models} to the same map-level skies, at the two survey sizes
of Section~\ref{sec:sims:survey} and with the prior of Table~\ref{tab:prior}. We show only the four
parameters that the counts constrain, namely $\sigma_8$, $\omc$, $\asz$ and $\alsz$. Since the linear
calibration is trained at 10 per cent of the sky only, it appears only in the coverage test and in the
$\sigma_8$ bias averaged over the prior. The size-split calibration appears only in
Section~\ref{sec:calib-size}. In each subsection we present the map-level posterior first, then the
catalogue models calibrated on maps, and finally the uncalibrated catalogue model, which serves as a
reference for the size of the finder's effects.

\subsection{Bias at the true parameters}\label{sec:results:bias}

We first apply each posterior to 100 map-level mocks at the true parameters and measure its bias
$\mathcal{B}$ (equation~\ref{eq:bias}), which Fig.~\ref{fig:bias} plots for each parameter and each
posterior and Table~\ref{tab:bias} lists. With 100 mocks, the errors of the mean are
0.04--0.11, so that biases of about 0.1--0.2 widths are resolved.

\begin{table}
\centering
\caption{Bias $\mathcal{B}$ of each posterior at the true parameters, in posterior widths, defined as the
mean over 100 map-level mocks of (posterior mean $-$ truth)/(posterior width) and given with its error of
the mean. The networks combined (Section~\ref{sec:inference:ensembles}) are, at 10 per cent of the sky, five
for the map-level posterior and one for each of the others, and at 40.7 per cent, four for the uncalibrated
model (the fifth training run did not finish within its time limit) and five for each of the others. The
same values are plotted in Fig.~\ref{fig:bias}.}
\label{tab:bias}
\small
\medmuskip=2mu
\setlength{\tabcolsep}{2pt}
\begin{tabular}{@{}lcccc@{}}
\toprule
 & $\sigma_8$ & $\omc$ & $\asz$ & $\alsz$ \\
\midrule
\multicolumn{5}{l}{\emph{10 per cent of the sky (41 tiles)}}\\
uncalibrated     & $+1.86\pm0.10$ & $-0.77\pm0.06$ & $-1.42\pm0.08$ & $-5.28\pm0.11$ \\
one-point        & $+0.13\pm0.09$ & $+0.09\pm0.06$ & $-0.06\pm0.08$ & $+0.50\pm0.08$ \\
noise-level      & $-0.07\pm0.10$ & $+0.11\pm0.06$ & $+0.09\pm0.08$ & $+0.62\pm0.09$ \\
map-level        & $-0.05\pm0.10$ & $+0.10\pm0.07$ & $-0.05\pm0.08$ & $+0.08\pm0.08$ \\
\midrule
\multicolumn{5}{l}{\emph{40.7 per cent of the sky (168 tiles)}}\\
uncalibrated     & $+2.85\pm0.07$ & $-1.19\pm0.04$ & $-1.74\pm0.06$ & $-8.30\pm0.07$ \\
one-point        & $+0.20\pm0.08$ & $+0.17\pm0.04$ & $-0.11\pm0.06$ & $+0.52\pm0.06$ \\
noise-level      & $+0.05\pm0.07$ & $+0.21\pm0.04$ & $-0.03\pm0.06$ & $+0.67\pm0.07$ \\
map-level        & $-0.11\pm0.08$ & $+0.10\pm0.04$ & $-0.05\pm0.06$ & $+0.03\pm0.06$ \\
\bottomrule
\end{tabular}
\end{table}

We find that the map-level posterior (blue) is within 0.12 widths of the truth in every parameter at both
survey sizes. Its largest departure in magnitude is $-0.11\pm0.08$ in $\sigma_8$ at 40.7 per cent of the
sky. Its most significant one, $+0.10\pm0.04$ widths in $\omc$ at the same size, amounts to 2.2 times its
error (from the unrounded values), and we return to it in Section~\ref{sec:discussion:omegac}. The bias
of a single network moves by about 0.1 widths from seed to seed (Appendix~\ref{app:robust}).

The two calibrated catalogue models (green and violet) are also within 0.21 widths of the truth in
$\sigma_8$, $\omc$ and $\asz$ at both survey sizes. However, they leave $\alsz$ high, by 0.50 and 0.52
widths with the one-point calibration and by 0.62 and 0.67 widths with the noise-level calibration, at 10
and 40.7 per cent of the sky, respectively. Because the one-point calibration uses the finder's noise
measured on map-level mocks, we discuss in Section~\ref{sec:discussion:calibrated} what this result implies
for a catalogue model with the noise measured on the maps.

Finally, the uncalibrated catalogue model (orange) serves as the reference, showing how large the finder's
effects are at goal depth when the catalogue model leaves them out. On our skies it is strongly biased, and
at 10 per cent of the sky it places $\sigma_8$ 1.86 widths high and $\alsz$ 5.28 widths low, with $\omc$ and
$\asz$ also low. Since the shifts in $\sigma_8$ and $\Omm$ partly cancel in
$S_8 \equiv \sigma_8(\Omm/0.3)^{1/2}$, the bias in $S_8$ is only 0.30 widths, again low.

From 10 to 40.7 per cent of the sky, the bias of the uncalibrated model in widths grows by a factor of about
1.5, to $+2.85$ in $\sigma_8$ and $-8.30$ in $\alsz$, and by a factor of 1.2, to $-1.74$, in $\asz$. This
growth does not come from the shift of the posterior mean, which hardly changes with area, going from
$+0.0115$ to $+0.0103$ in $\sigma_8$ and from $-0.086$ to $-0.084$ in $\alsz$. The posteriors, in
contrast, become 1.3--1.7 times narrower, by factors of 1.7 in $\sigma_8$, 1.6 in $\alsz$, 1.4 in $\asz$
and 1.3 in $\omc$, which is less than the factor of 2 expected from four times the area if the counts alone
set the widths. It is this narrowing that makes the bias in widths grow. Because the two survey sizes
share their first 41 tiles (Section~\ref{sec:sims:survey}), these biases are not independent. All of them
hold for our set-up, that is, for goal-depth forecast noise, a noise table for the instrument and the CMB
only, and a flat prior on the slope $\alsz$ (Section~\ref{sec:discussion:scope}).

\subsection{The posteriors of one mock}\label{sec:results:corner}

\begin{figure}
\centering
\includegraphics[width=\columnwidth]{figures/corner.png}
\caption{Posteriors of the first map-level mock of the 41-tile survey (10 per cent of the sky) in
$\sigma_8$, $\Omm$, $\asz$ and $\alsz$, from three posteriors of one network each: the map-level posterior
(blue), the noise-level calibration (violet) and, as a reference, the uncalibrated catalogue model
(orange). The off-diagonal panels show the 68 and 95 per cent contours of the two-dimensional marginals
(thick and thin lines), and hence how the parameters trade off against each other, and the diagonal panels
compare the one-dimensional marginals with the truth. Black dots and dotted lines mark the true values.
$\Omm$ is derived from $\omc$, $\omb$ and $h$.}
\label{fig:corner}
\end{figure}

\begin{figure}
\centering
\includegraphics[width=\columnwidth]{figures/corner_40.png}
\caption{As Fig.~\ref{fig:corner}, for the first map-level mock of the 168-tile survey (40.7 per cent of
the sky), with four posteriors: the map-level posterior (blue, five networks pooled), the noise-level
calibration (violet, five networks pooled), the one-point calibration (green, one network) and, as a
reference, the uncalibrated catalogue model (orange, one network). Contours are at 68 and 95 per cent, the
diagonal panels show the one-dimensional marginals, and black marks the true values.}
\label{fig:corner40}
\end{figure}

\begin{figure*}
\centering
\includegraphics[width=\textwidth]{figures/coverage.png}
\caption{Coverage on 500 map-level skies drawn from the prior, for $\sigma_8$, $\omc$, $\asz$ and $\alsz$
(columns), at 10 per cent (top row) and 40.7 per cent of the sky (bottom row). The horizontal axis is the
credible level $\gamma$ of the central interval, and the vertical axis is $\cov(\gamma)$, the fraction of
skies whose true value lies in that interval, obtained from the rank of the truth among 4000 posterior
samples. Green shows the one-point calibration, dashed grey the linear calibration (10 per cent only),
violet the noise-level calibration and blue the map-level posterior. At 10 per cent of the sky each
catalogue-level posterior is one network and the map-level posterior is five, while at 40.7 per cent the
one-point calibration is one network and the noise-level calibration and map-level posterior are five
each. The dotted diagonal $\cov(\gamma)=\gamma$ is that of a posterior with correct coverage, and a curve
below it means intervals that are too narrow or off-centre in part of the prior. The uncalibrated model is
not shown, because on map-level skies most of its samples fall outside the prior and are rejected, which
makes sampling 500 skies impractical. Values at $\gamma = 0.5$ and 0.9 are in Table~\ref{tab:coverage}.}
\label{fig:coverage}
\end{figure*}

Since the biases of Table~\ref{tab:bias} are averages over 100 mocks, we show in Figs~\ref{fig:corner}
and~\ref{fig:corner40} what they mean for a single sky, using the posterior of the first mock in the four
parameters that the counts constrain best. In place of $\omc$, these corner plots show the derived $\Omm$,
the density that enters the mass function. At 10 per cent of the sky each of these posteriors, including
the map-level one, comes from a single network and can therefore differ slightly from the ensembles of
Table~\ref{tab:bias}.

At 10 per cent of the sky (Fig.~\ref{fig:corner}), we find that the map-level posterior contains the truth, and the
noise-level calibration almost coincides with it. In contrast, the posterior trained on the uncalibrated
model misses the truth, placing $\sigma_8$ about two widths high, $\alsz$ several widths low, and $\Omm$
and $\asz$ low. These are the directions that allow it to explain the deficit of bright clusters on the
maps with a shallower mass slope. Despite these shifts, the three posteriors have similar widths, which
shows that on this data vector map-level inference is not more precise than catalogue-level inference.

At 40.7 per cent of the sky (Fig.~\ref{fig:corner40}), the posteriors are narrower. The map-level
posterior gives $\sigma_8 = 0.812\pm0.003$ and $\alsz = 1.120\pm0.010$, against true values of 0.811 and
1.123, and the two calibrated models agree with it and contain the truth. On a single mock at the true
parameters, we therefore cannot tell the calibrated catalogue models from map-level inference. They differ
from it, however, in two ways that one mock cannot reveal, namely the shift of $\alsz$ by 0.5--0.7 widths
over 100 mocks (Table~\ref{tab:bias}) and their coverage away from the true parameters, which we test in
Section~\ref{sec:results:coverage}. The bias of the uncalibrated model is larger at this area, and on this mock
it gives $\sigma_8 = 0.8231\pm0.0036$, 3.4 widths from the truth, and $\alsz = 1.034\pm0.011$,
8.3 widths from it.

\begin{figure}
\centering
\includegraphics[width=\columnwidth]{figures/bias_thirds.png}
\caption{Bias of $\sigma_8$ (posterior mean $-$ truth, in posterior widths) on the 500 map-level skies
drawn from the prior at 10 per cent of the sky, averaged in thirds (low, middle, high) of each sky's true
$\omc$ (left), $\asz$ (centre) and $\alsz$ (right). In each panel the thirds run from the lowest to the
highest true value of one parameter, so that a posterior unbiased everywhere would give a flat row of
points at zero. Green shows the one-point calibration and violet the noise-level calibration (one network
each). The dotted line marks zero bias.}
\label{fig:thirds}
\end{figure}

\subsection{Coverage over the prior}\label{sec:results:coverage}

\begin{table}
\centering
\caption{Coverage $\cov(\gamma)$ at $\gamma = 0.5$ and 0.9, given as $\cov(0.5)/\cov(0.9)$, that is, the
fraction of 500 map-level skies drawn from the prior whose true value lies in the central 50 and 90 per
cent credible intervals. A posterior with correct coverage has $0.5/0.9$, and with 500 skies the binomial
error is 0.022 at $\gamma=0.5$ and 0.013 at $\gamma=0.9$. The number of networks combined is given in
brackets, and the full curves are shown in Fig.~\ref{fig:coverage}.}
\label{tab:coverage}
\small
\setlength{\tabcolsep}{3.5pt}
\begin{tabular}{lcccc}
\toprule
 & $\sigma_8$ & $\omc$ & $\asz$ & $\alsz$ \\
\midrule
\multicolumn{5}{l}{\emph{10 per cent of the sky (41 tiles)}}\\
one-point (1)    & 0.32/0.65 & 0.40/0.72 & 0.41/0.82 & 0.39/0.80 \\
linear (1)       & 0.42/0.84 & 0.46/0.86 & 0.47/0.88 & 0.41/0.79 \\
noise-level (1)  & 0.46/0.90 & 0.47/0.89 & 0.48/0.88 & 0.42/0.78 \\
map-level (5)    & 0.46/0.90 & 0.47/0.87 & 0.49/0.89 & 0.44/0.86 \\
\midrule
\multicolumn{5}{l}{\emph{40.7 per cent of the sky (168 tiles)}}\\
one-point (1)    & 0.30/0.59 & 0.35/0.68 & 0.40/0.79 & 0.37/0.77 \\
noise-level (5)  & 0.50/0.86 & 0.51/0.90 & 0.52/0.90 & 0.35/0.74 \\
map-level (5)    & 0.51/0.90 & 0.47/0.88 & 0.48/0.91 & 0.46/0.88 \\
\bottomrule
\end{tabular}
\end{table}

A small bias at the true parameters can hide a posterior that fails elsewhere in the prior, and we
therefore also test the coverage over the whole prior. We draw 500 skies from the prior, run the map-level
simulator and the finder on each, and count how often the central credible interval of level $\gamma$
contains the true value (equation~\ref{eq:coverage}). Fig.~\ref{fig:coverage} plots this fraction,
$\cov(\gamma)$, against $\gamma$, and Table~\ref{tab:coverage} gives $\cov(0.5)$ and $\cov(0.9)$.

The map-level posterior (blue) lies close to the diagonal for all four parameters at both survey sizes,
with 90 per cent coverage between 0.86 and 0.91 and, at 40.7 per cent of the sky, 50 per cent coverage
between 0.46 and 0.51. At 10 per cent of the sky we find its $\alsz$ intervals slightly too narrow, with
a coverage of 0.86, about three binomial errors below 0.9, and its $\omc$ intervals likewise (0.87). At
40.7 per cent, however, both have close to correct coverage (0.88). We show in
Appendix~\ref{app:robust} that the value of 0.86 depends neither on the training seed nor on the number of
simulations.

In contrast, the one-point calibration (green) falls well below the diagonal for $\sigma_8$ and $\omc$. Its 90 per cent intervals contain the true $\sigma_8$ for only 0.65 of the skies at 10 per cent of the
sky and for 0.59 at 40.7 per cent. The intervals are off-centre over much of the prior
(Section~\ref{sec:results:thirds}), and this becomes more apparent as the survey grows and the posteriors
narrow. A calibration that follows the parameters brings $\sigma_8$, $\omc$ and $\asz$ close to the
diagonal. Their 90 per cent coverage is 0.84--0.88 for the linear calibration (dashed, 10 per cent of the
sky only) and 0.86--0.90 for the noise-level calibration (violet) at both survey sizes. Neither calibration, however, corrects $\alsz$. Its 90 per cent coverage is 0.79 with the linear and 0.78 with
the noise-level calibration at 10 per cent of the sky, and 0.74 with the noise-level calibration at 40.7
per cent. We do not apply this test to the uncalibrated model, because on map-level skies most of its samples
fall outside the prior and are rejected.

\subsection{Where the calibrated models fail}\label{sec:results:thirds}

To see where in the prior the calibrated models fail, we split the 500 prior skies at 10 per cent of the
sky into thirds of each sky's true $\omc$, $\asz$ and $\alsz$, and average the bias of $\sigma_8$ in each
third (Fig.~\ref{fig:thirds}). Because each third averages over the other parameters, no third coincides
with the single point at which the one-point calibration is fitted.

The bias of the one-point calibration (green) depends strongly on where the sky lies in the prior. It is
close to zero only in the lowest thirds of $\omc$ ($+0.02$) and $\asz$ ($-0.09$), where the total tSZ
signal on the sky is smallest. It increases with the tSZ signal, reaching $+0.50$ and $+2.64$ widths
in the middle and highest thirds of $\omc$ and $+1.00$ and $+2.26$ widths in those of $\asz$. The reason is
that, where the true $\omc$ or $\asz$ is high, the sky contains more tSZ signal from faint clusters than at the calibration point. The finder's noise is then higher, fewer bright clusters are detected, and the
posterior interprets the deficit as a higher $\sigma_8$. Since a lower slope $\alsz$ also means more signal
from faint clusters, the trend is reversed in $\alsz$, with $+1.70$ widths in its lowest third, $+0.85$ in
the middle and $+0.63$ in the highest. Averaged over the whole prior, the $\sigma_8$ bias of the one-point
calibration is $+1.06$ widths.

With the noise-level calibration (violet), the model follows these changes and the trends largely
disappear. Its bias lies between $+0.04$ and $+0.26$ widths in every third except the highest thirds of
$\omc$ and $\asz$, which contain the skies with the most tSZ signal and where it reaches $+0.47$ and
$+0.58$, respectively. Over the whole prior, the $\sigma_8$ bias is $+0.32$ widths for the linear
calibration and $+0.23$ for the noise-level calibration. This residual is largest where the tSZ signal is
highest and the finder's errors have the heaviest upward tail, which a catalogue model with Gaussian
scatter does not include (Section~\ref{sec:finder_tail}). Since Fig.~\ref{fig:thirds} has neither a
map-level reference nor error bars, however, it does not show how much of the residual is due to the
calibration itself. Part of it may also come from the edges of the prior, because averaging in thirds of a
parameter correlated with $\sigma_8$ can move the mean posterior offset away from zero even for a
posterior with correct coverage.
